\documentclass[a4paper,fleqn,review,12pt]{cas-sc}

\usepackage[numbers,sort&compress]{natbib}
\usepackage{amsmath,amssymb,bm}
\usepackage{graphicx}
\usepackage{booktabs}
\usepackage{xcolor}
\usepackage{xurl}
\usepackage{lineno}
\usepackage[section]{placeins}

\graphicspath{{submit_figure/}}
\hypersetup{hidelinks}
\setcounter{secnumdepth}{3}
\setlength{\emergencystretch}{2em}
\allowdisplaybreaks
\renewcommand{\floatpagefraction}{.65}
\AddToHook{env/thebibliography/begin}{\interlinepenalty=10000}
\numberwithin{equation}{section}

% C5 text is current; non-DOI revision wrappers are removed while DOI placeholders remain red.
\newcommand{\rev}[1]{\textcolor{red}{#1}}
\newcommand{\todo}[1]{\textcolor{red}{\textbf{[TODO: #1]}}}

\begin{document}
\let\WriteBookmarks\relax
\def\floatpagepagefraction{1}
\def\textpagefraction{.001}

\title[mode=title]{On Craig--Bampton Model Reduction for Multi-Degree-of-Freedom Coupled Real-Time Hybrid Simulation}

\author{Ning Li}
\author{Yu Liang}
% Affiliations, e-mail addresses, and the corresponding author were not supplied in temp.tex.

\begin{abstract}
Limited actuation channels motivate the condensation of interface degrees of freedom (DOFs) in multi-degree-of-freedom real-time hybrid simulation (RTHS), changing both structural response and delay-dependent stability.
This study develops a joint accuracy and stability formulation for Guyan and Craig--Bampton condensation under a fixed actuation constraint.
A common projection framework enforces compatibility at the retained interface DOFs, and channel-specific delay operators identify how reconstructed responses enter the feedback loop.
The discrete-time characteristic equation incorporates the CR integration algorithm and linear quadratic regulator feedback, with closed-loop poles determining admissible actuator-delay combinations.
Two substructure divisions of a linear three-storey frame are examined through virtual RTHS with two actuation channels.
Delay-free \rev{full-frame} accuracy comparisons give a maximum relative error of \rev{0.284\%} in the first two natural frequencies for Craig--Bampton condensation.
Its first-storey displacement normalized root mean square error \rev{has a maximum of 0.138\%} for the earthquake record and all evaluated chirp bands when normalized by the full-record response range.
The Guyan model that condenses the middle-storey horizontal DOF produces corresponding maxima of \rev{30.574\% and 26.858\%}.
With independent constant actuator delays, Craig--Bampton condensation yields a larger stable delay region than Guyan condensation in both divisions.
\rev{Expanded modal shapes and reconstructed middle-storey responses identify the response components affected by condensation.}
The results connect the retained internal modal dynamics to response fidelity and \rev{the delay tolerance of each model--regulator pair}, providing a basis for reduced-model selection under limited actuation capacity.

\end{abstract}

\begin{keywords}
	\raggedright
	Real-time hybrid simulation
	\sep Craig--Bampton condensation
	\sep Guyan condensation
	\sep Actuation constraints
	\sep Multiple actuator delays
	\sep Delay-dependent stability
\end{keywords}

\shortauthors{N. Li and Y. Liang}
\shorttitle{\mbox{}}

\maketitle
\linenumbers

\section{Introduction} \label{sec1}

Reliable assessment of structural performance requires experimental characterization of components whose nonlinear or rate-dependent behaviour is difficult to reproduce numerically. Testing the entire structure, however, is restricted by specimen size, loading capacity and cost. Real-time hybrid simulation (RTHS) addresses this difficulty by coupling a physical substructure containing the components of interest with a numerical representation of the remaining structure, while preserving the time scale of the imposed loading \cite{Nakashima1992}. This combination supports the investigation of multi-directional building response under natural hazards \cite{AlSubMulti2024}, seismic soil--foundation--structure interaction \cite{MalikMulti2025}, and wind-induced aeroelastic response \cite{MalikAero2026}. Its effectiveness depends on reproducing the interaction between the substructures with sufficient spatial resolution, tracking accuracy and real-time stability.

Accurate interface loading requires the actuators to follow the numerical commands despite their dynamics and interaction with the specimen. Model-based prediction and compensation address these effects by incorporating information about the loading system and structural response \cite{Carrion2007}. For multi-degree-of-freedom (MDOF) systems, mechanical coupling and unequal actuator responses make this a multi-channel control problem, as represented by the multi-axial RTHS benchmark \cite{Condori2023}. Recent approaches update compensation parameters during loading through robust decentralized adaptation and conditional adaptive time-series methods \cite{Quiroz2024,Aguila2024}. Sliding-mode compensation and two-stage $H_{\infty}$ control explicitly address coupled actuator dynamics and robustness \cite{Shangguan2024,Xie2025}, while discrete feedforward control improves command tracking within the sampled-data implementation \cite{Calayir2025}. Nonlinear parameter estimation and unscented Kalman filtering provide another route to updating the compensation model \cite{Ruiz2024,WangUKF2024}. Data-driven nonlinear autoregressive models and long short-term memory networks further represent actuator behaviour and improve response prediction \cite{XuNARX2024,Lai2025}. These methods improve the execution of commands assigned to existing actuator channels. A separate limitation arises when the physical substructure contains more interface degrees of freedom (DOFs) than the available actuators can impose independently. Improving tracking on the actuated DOFs does not supply the missing loading conditions on the remaining interface.

The real-time computational requirement introduces a second constraint. Explicit integration algorithms reduce the cost of advancing structural states \cite{ChenCR2008}, and recent dissipative formulations support multi-directional simulation of large nonlinear systems \cite{AlSubCD2024}. Multi-rate execution separates the numerical and control update rates, while concurrent computing improves the scheduling of multi-axial simulation tasks \cite{TaoMulti2024,Sudvarg2024}. Neural-network representations reduce the online cost of nonlinear numerical substructures and soil--foundation interaction models \cite{ChenNN2024,AlSubNN2024}. Online cyber-physical neural networks also use measurements from a tested device to represent devices at other locations, reducing the number of components that must be physically implemented \cite{MalikNN2025}. An alternative is to replace the online interaction by successive full-record numerical and physical runs. Offline iterative methods have been developed for seismic fluid--structure interaction \cite{HuOffline2024}, with signal mapping and model updating used to improve convergence \cite{HuIntegrated2025,HuSequence2025}. Neural-network-enhanced time-history iteration provides a further means of accelerating this process \cite{Peng2025}. The corresponding convergence analysis shows why intermediate iteration errors, rather than only the final converged response, require attention \cite{HuangGreen2026}. These developments address computational cost, replication of component behaviour or the temporal organization of the test. For an online RTHS with a prescribed physical substructure, the spatial approximation introduced by a limited number of independently imposed interface DOFs remains to be assessed.

Residual actuator delay must also be evaluated in the assembled feedback loop. Discrete-time stability formulations account for the combined effects of numerical integration and actuator delay in single-DOF systems \cite{ChenStability2008}, and extend the analysis to MDOF systems through discrete root loci and pole-based formulations \cite{Zhu2015,Li2021}. Lyapunov--Krasovskii formulations distinguish constant-delay stability from stability under time-varying delay \cite{Huang2019}. More recent analysis of inerter-type experimental substructures demonstrates the importance of the physical feedback mechanism when constructing a delay-dependent model \cite{TaoInerter2025}. Local control assessment using instantaneous frequency and compensation based on amplitude and phase errors additionally distinguish tracking quality from a single equivalent delay value \cite{XuInst2024,XuAmp2025}. The increasing dimensionality and coupling of experimental boundaries have likewise motivated a complexity-based classification of boundary-control problems \cite{Zhan2026}. For a condensed online system, however, the delayed response components are determined jointly by the actuator assignment and the reduction basis. A stability analysis formulated for a specified set of delayed DOFs does not, by itself, quantify the consequences of replacing additional interface DOFs by reconstructed responses.

Condensation provides a direct representation of this spatial approximation. Guyan condensation expresses the eliminated DOFs through the static deformation induced by the retained DOFs \cite{Guyan1965}. Craig--Bampton condensation enriches this representation with selected fixed-interface modes \cite{Craig1968}. Classical Craig--Bampton assembly retains all interface DOFs, and reducing the interface requires an additional approximation \cite{Krattiger2019}. Model reduction is already established in hybrid simulation. Component-mode synthesis has been used to match the relevant structural frequency range to actuator bandwidth \cite{Miraglia2020}, dynamic condensation has reduced the online computational burden of RTHS \cite{Mucha2023}, and proper orthogonal decomposition has enabled the solution of large numerical substructures \cite{ZhangPOD2024}. The present issue is the joint effect of condensation and independent actuator delays when the retained physical DOFs are selected by the available loading channels. The Guyan reconstruction is entirely determined by the retained physical DOFs, whereas the Craig--Bampton reconstruction also contains numerically evolved modal DOFs. Consequently, the two representations distribute structural response differently between explicitly delayed actuator channels and the remaining numerical states. Preserving natural frequencies or retained-DOF mode shapes alone cannot establish the accuracy and stability of that interaction. The reduction basis, reduced-interface coupling and channel delays must be examined together under the specified actuation constraint.

This paper establishes a joint accuracy and delay-dependent stability formulation for actuation-constrained MDOF RTHS. Guyan and Craig--Bampton condensations are expressed within a common projection framework, with compatibility enforced at the retained interface DOFs. The reconstruction is then combined with channel-specific delay operators to identify the response components transmitted through each actuator. A discrete-time closed-loop characteristic equation incorporates the CR integration algorithm and a linear quadratic regulator, and admissible delay combinations are determined from the closed-loop poles. Modal and time-domain errors assess the structural approximation, while a modal redistribution ratio measures the change in low-order modal content associated with the actuator-driven DOFs. Together, these quantities connect the choice of condensation basis and substructure division to response fidelity and delay tolerance, providing a basis for selecting reduced models under a fixed number of actuation channels.


\section{Structural modeling and condensation formulations} \label{sec2}

In a multi-degree-of-freedom (MDOF) real-time hybrid simulation (RTHS), the physical substructure often contains more interface coordinates than the transfer system can impose independently. Model reduction is adopted to express the coordinates without independent actuation through a reduced set of experimentally imposed coordinates.
In this section, the structural modeling framework and the condensation formulations are established. The full-order equation of motion is decomposed into the contributions of the numerical and physical substructures. Guyan static condensation and Craig--Bampton dynamic condensation are then formulated on a common partition of retained and condensed coordinates, and the two methods differ only in the construction of the transformation matrix.

\subsection{Equations of motion and substructure partitioning}
The equation of motion of the full-order structure under a ground acceleration is expressed as

\vspace{-3.5em}
\begin{equation}
	\mathbf{M}\ddot{\bm{x}}(t)+\mathbf{C}\dot{\bm{x}}(t)+\mathbf{K}\bm{x}(t)
	=-\mathbf{M}\bm{\Gamma}a_{\mathrm{g}}(t),
	\label{eq:full-order-eom}
\end{equation}
\vspace{-3.5em}

\noindent where \(\mathbf{M}\), \(\mathbf{C}\), and \(\mathbf{K}\) are the mass, damping, and stiffness matrices of the full-order structure, respectively, \(\bm{x}(t)\) is the generalized displacement vector of the full-order structural DOFs measured relative to the ground, \(\bm{\Gamma}\) is the influence vector, and \(a_{\mathrm{g}}(t)\) is the ground acceleration.

After the structure is divided into numerical and physical substructures, the corresponding equations are written as

\vspace{-3.5em}
\begin{subequations}
	\label{eq:substructure-eom}
	\begin{align}
		\mathbf{M}_{\mathrm{n}}\ddot{\bm{x}}_{\mathrm{n}}
		+\mathbf{C}_{\mathrm{n}}\dot{\bm{x}}_{\mathrm{n}}
		+\mathbf{K}_{\mathrm{n}}\bm{x}_{\mathrm{n}}
		=\bm{f}_{\mathrm{n}}{}+\bm{\lambda}_{\mathrm{n}},
		\label{eq:numerical-eom}\\
		\mathbf{M}_{\mathrm{e}}\ddot{\bm{x}}_{\mathrm{e}}
		+\mathbf{C}_{\mathrm{e}}\dot{\bm{x}}_{\mathrm{e}}
		+\mathbf{K}_{\mathrm{e}}\bm{x}_{\mathrm{e}}
		=\bm{f}_{\mathrm{e}}{}+\bm{\lambda}_{\mathrm{e}},
		\label{eq:physical-eom}
	\end{align}
\end{subequations}
\vspace{-3.5em}

\noindent where \(\bm{x}_{\mathrm{n}}\) and \(\bm{x}_{\mathrm{e}}\) are the DOFs of the numerical and physical substructures, respectively, \(\bm{f}_{\mathrm{n}}\) and \(\bm{f}_{\mathrm{e}}\) are the external load vectors carrying the share of the seismic inertia load of Eq. \eqref{eq:full-order-eom} taken by each substructure, and \(\bm{\lambda}_{\mathrm{n}}\) and \(\bm{\lambda}_{\mathrm{e}}\) are the interface force vectors acting on the two substructures. \rev{The interface reactions balance after the local force vectors are mapped to the same interface coordinates.} The partition assigns one part of the structure to numerical solution and the other to physical implementation, and it does not dynamically decouple the two substructures.

For a condensation operation, the DOFs of a substructure are arranged into a retained set \(\bm{u}_{\mathrm{r}}\) and a condensed set \(\bm{u}_{\mathrm{c}}\), and the equation of motion takes the form

\vspace{-3.5em}
\begin{equation}
	\begin{bmatrix}
		\mathbf{M}_{\mathrm{rr}} & \mathbf{M}_{\mathrm{rc}}\\
		\mathbf{M}_{\mathrm{cr}} & \mathbf{M}_{\mathrm{cc}}
	\end{bmatrix}
	\begin{bmatrix}
		\ddot{\bm{u}}_{\mathrm{r}}\\
		\ddot{\bm{u}}_{\mathrm{c}}
	\end{bmatrix}
	+
	\begin{bmatrix}
		\mathbf{C}_{\mathrm{rr}} & \mathbf{C}_{\mathrm{rc}}\\
		\mathbf{C}_{\mathrm{cr}} & \mathbf{C}_{\mathrm{cc}}
	\end{bmatrix}
	\begin{bmatrix}
		\dot{\bm{u}}_{\mathrm{r}}\\
		\dot{\bm{u}}_{\mathrm{c}}
	\end{bmatrix}
	+
	\begin{bmatrix}
		\mathbf{K}_{\mathrm{rr}} & \mathbf{K}_{\mathrm{rc}}\\
		\mathbf{K}_{\mathrm{cr}} & \mathbf{K}_{\mathrm{cc}}
	\end{bmatrix}
	\begin{bmatrix}
		\bm{u}_{\mathrm{r}}\\
		\bm{u}_{\mathrm{c}}
	\end{bmatrix}
	=
	\begin{bmatrix}
		\bm{f}_{\mathrm{r}}\\
		\bm{f}_{\mathrm{c}}
	\end{bmatrix},
	\label{eq:partitioned-local-eom}
\end{equation}
\vspace{-3.5em}

\noindent where \(\bm{f}_{\mathrm{r}}\) and \(\bm{f}_{\mathrm{c}}\) are the generalized forces acting on the retained and condensed coordinates, which collect the seismic load carried by the substructure together with the interface force.

The condensed coordinates are expressed through the retained coordinates by a transformation matrix, and the substructure matrices are projected onto the reduced subspace by a congruent transformation,

\vspace{-3.5em}
\begin{equation}
	\begin{gathered}
		\bm{u}(t)=\mathbf{T}\bm{q}(t),\\
		\mathbf{M}_{\mathrm{red}}=\mathbf{T}^{\mathsf{T}}\mathbf{M}\mathbf{T},\quad
		\mathbf{C}_{\mathrm{red}}=\mathbf{T}^{\mathsf{T}}\mathbf{C}\mathbf{T},\quad
		\mathbf{K}_{\mathrm{red}}=\mathbf{T}^{\mathsf{T}}\mathbf{K}\mathbf{T},\quad
		\bm{f}_{\mathrm{red}}=\mathbf{T}^{\mathsf{T}}\bm{f},
	\end{gathered}
	\label{eq:general-projection}
\end{equation}
\vspace{-3.5em}

\noindent where \(\bm{q}(t)\) is the reduced generalized-coordinate vector and \(\mathbf{T}\) is the transformation matrix. Guyan and Craig--Bampton condensation share Eq. \eqref{eq:general-projection} and differ only in the construction of \(\mathbf{T}\).

The two substructures are condensed separately and are coupled at the interface coordinates retained by both reduced models, \rev{where the projected interface forces balance. The reduced coordinates of substructure \(s\) satisfy \(\bm q_s=\mathbf A_s\bm q\), where \(\mathbf A_s\) places its local coordinates in the assembled coordinate vector.} The assembled reduced model is governed by

\vspace{-3.5em}
\begin{equation}
	\mathbf{M}_{\mathrm{red}}\ddot{\bm{q}}(t)
	+\mathbf{C}_{\mathrm{red}}\dot{\bm{q}}(t)
	+\mathbf{K}_{\mathrm{red}}\bm{q}(t)
	=\bm{f}_{\mathrm{red}}(t),
	\label{eq:assembled-reduced}
\end{equation}
\vspace{-3.5em}

\noindent where \(\bm{q}(t)\) collects the generalized coordinates of the two reduced substructures, with the shared interface coordinates counted once. \rev{The reduced substructure matrices assemble as \(\mathbf M_{\mathrm{red}}=\sum_s\mathbf A_s^{\mathsf T}\mathbf M_{s,\mathrm{red}}\mathbf A_s\), with the same operation for damping, stiffness and external loads, where \(s\in\{\mathrm n,\mathrm e\}\) identifies the numerical and physical parts. Interface forces satisfy \(\sum_s\mathbf A_s^{\mathsf T}\mathbf T_s^{\mathsf T}\bm\lambda_s=\bm0\) in this coupling space. Compatibility is imposed at the shared retained coordinates. At a geometric interface coordinate outside this set, each substructure has its own recovered displacement, so the two recovered values need not coincide.} The same restriction applies to the test itself. Once the interface carries more coordinates than there are actuation channels, the specimen responds on the unactuated interface coordinates according to its own dynamics, and \rev{the specified two-channel loading system does not impose an independent numerical command there.}

\subsection{Guyan static condensation} \label{sec21}

Guyan condensation derives the condensed coordinates from static equilibrium \cite{Guyan1965}. The reduction basis is built from the static response of the condensed coordinates. Their inertia, their damping and the load acting on them are therefore dropped from the lower block of Eq. \eqref{eq:partitioned-local-eom}, which leaves \(\mathbf{K}_{\mathrm{cr}}\bm{u}_{\mathrm{r}}+\mathbf{K}_{\mathrm{cc}}\bm{u}_{\mathrm{c}}=\bm{0}\). The condensed coordinates are approximated by

\vspace{-3.5em}
\begin{equation}
	\bm{u}_{\mathrm{c}}=\bm{\Psi}_{\mathrm{G}}\bm{u}_{\mathrm{r}},
	\qquad
	\bm{\Psi}_{\mathrm{G}}=-\mathbf{K}_{\mathrm{cc}}^{-1}\mathbf{K}_{\mathrm{cr}},
	\label{eq:guyan-coordinate-relation}
\end{equation}
\vspace{-3.5em}

\noindent where \(\bm{\Psi}_{\mathrm{G}}\) is the static condensation matrix, which requires \(\mathbf{K}_{\mathrm{cc}}\) to be non-singular.
The Guyan transformation is

\vspace{-3.5em}
\begin{equation}
	\bm{u}=\mathbf{T}_{\mathrm{G}}\bm{q}_{\mathrm{G}},
	\qquad
	\mathbf{T}_{\mathrm{G}}=
	\begin{bmatrix}
		\mathbf{I}\\
		\bm{\Psi}_{\mathrm{G}}
	\end{bmatrix},
	\qquad
	\bm{q}_{\mathrm{G}}=\bm{u}_{\mathrm{r}},
	\label{eq:guyan-transformation}
\end{equation}
\vspace{-3.5em}

\noindent where \(\bm{q}_{\mathrm{G}}\) is the reduced coordinate vector, which contains the retained coordinates alone. The reduced matrices follow from Eq. \eqref{eq:general-projection} with \(\mathbf{T}=\mathbf{T}_{\mathrm{G}}\), and the response of the condensed coordinates is recovered from Eq. \eqref{eq:guyan-coordinate-relation}.

The loads on the condensed coordinates, including their share of the seismic inertia load and the interface force, are projected through \(\mathbf{T}_{\mathrm{G}}^{\mathsf{T}}\bm{f}\), and they no longer act on independent coordinates. Guyan condensation is accordingly most accurate when the condensed coordinates follow the retained coordinates through quasi-static deformation, and it loses accuracy once the inertia of the condensed coordinates contributes to the response.

\subsection{Craig--Bampton dynamic condensation}
\label{sec:cb-condensation}

Craig--Bampton condensation keeps the retained coordinates explicit and represents the condensed coordinates through constraint modes and selected fixed-interface modes \rev{\cite{Craig1968}}. The experimentally imposed coordinates form the retained set of the physical substructure, and internal master coordinates are additionally retained in the numerical substructure. An interface coordinate driven by no actuator is therefore condensed, whereas the classical partition keeps every interface coordinate explicit \cite{Krattiger2019}.

The dynamic modes are obtained with the retained coordinates constrained, and satisfy

\vspace{-3.5em}
\begin{equation}
	\mathbf{K}_{\mathrm{cc}}\bm{\Phi}
	=\mathbf{M}_{\mathrm{cc}}\bm{\Phi}\bm{\Omega}^{2},
	\qquad
	\bm{\Phi}^{\mathsf{T}}\mathbf{M}_{\mathrm{cc}}\bm{\Phi}=\mathbf{I},
	\label{eq:fixed-interface-modes}
\end{equation}
\vspace{-3.5em}

\noindent where the columns of \(\bm{\Phi}\) are the fixed-interface modes of the partition and \(\bm{\Omega}^{2}\) is the diagonal matrix of the corresponding eigenvalues.

The constraint modes describe the static deformation of the condensed coordinates produced by unit displacements of the retained coordinates, which is the relation \(\bm{\Psi}_{\mathrm{G}}\) of Eq. \eqref{eq:guyan-coordinate-relation}. The Craig--Bampton transformation is therefore

\vspace{-3.5em}
\begin{equation}
	\begin{bmatrix}
		\bm{u}_{\mathrm{r}}\\
		\bm{u}_{\mathrm{c}}
	\end{bmatrix}
	=
	\begin{bmatrix}
		\mathbf{I} & \mathbf{0}\\
		\bm{\Psi}_{\mathrm{G}} & \bm{\Phi}_{d}
	\end{bmatrix}
	\begin{bmatrix}
		\bm{u}_{\mathrm{r}}\\
		\bm{\eta}
	\end{bmatrix}
	=
	\mathbf{T}_{\mathrm{CB}}\bm{q}_{\mathrm{CB}},
	\label{eq:cb-transformation}
\end{equation}
\vspace{-3.5em}

\noindent where \(\bm{\Phi}_{d}\) collects the \(d\) fixed-interface modes of lowest frequency, which are the modes retained in the reduced model, and \(\bm{\eta}\) contains the corresponding modal coordinates. The reduced matrices follow from Eq. \eqref{eq:general-projection} with \(\mathbf{T}=\mathbf{T}_{\mathrm{CB}}\).

The response of the condensed coordinates is approximated as

\vspace{-3.5em}
\begin{equation}
	\bm{u}_{\mathrm{c}}(t)=
	\bm{\Psi}_{\mathrm{G}}\bm{u}_{\mathrm{r}}(t)
	+\bm{\Phi}_{d}\bm{\eta}(t).
	\label{eq:cb-internal-recovery}
\end{equation}
\vspace{-3.5em}

The first term is the static contribution of Guyan condensation, and the second term retains selected internal vibration patterns, which is the only difference between the two transformations.


\section{Delay-dependent stability formulation} \label{sec3}

In an MDOF RTHS, each actuated coordinate of the physical substructure is driven by a separate actuator, and each actuator introduces its own response delay. Condensation expresses the condensed coordinates through the retained coordinates, and the delays acting on the retained coordinates are therefore carried by the reconstructed response of the condensed coordinates. In this section, the delay-dependent stability formulation of the condensed RTHS system is established. The actuator delays are represented as integer multiples of the integration time step, and their transmission through the Guyan and Craig--Bampton transformations is identified. The condensed substructures are then assembled into a discrete-time closed-loop model with an LQR controller, and the stability is assessed from the modulus of the dominant pole of the closed loop.

\subsection{Multiple-delay representation and delay propagation through condensation}
\label{sec:multiple-delay}

The delay of an RTHS system originates from the numerical integration, the data exchange between the numerical and physical substructures, the conversion between digital and analogue signals, and the response of the servo-hydraulic actuator \cite{Carrion2007}. The actuator response delay is retained in the present formulation and the other three sources are excluded, which isolates the interaction between the delay and the model reduction.

The actuator delay is taken as an integer multiple of the integration time step \cite{Zhu2015}. The delay then becomes a power of \(z\) after the \(z\) transform, which avoids the transcendental term produced by a continuous-time delay and admits several delay channels at the same time. The delay of the \(\ell\)th actuator and the displacement that this actuator imposes on the specimen are

\vspace{-3.5em}
\begin{equation}
	\tau_{\ell}=n_{\ell}\Delta t,
	\qquad
	\hat{u}_{\mathrm{r},\ell}(t)=u_{\mathrm{r},\ell}(t-\tau_{\ell}),
	\qquad
	\ell=1,\dots,n_{\mathrm{a}},
	\label{eq:integer-delay}
\end{equation}
\vspace{-3.5em}

\noindent where \(\Delta t\) is the integration time step, \(n_{\ell}\) is a non-negative integer, with \(n_{\ell}=0\) denoting a delay-free channel, \(n_{\mathrm{a}}\) is the number of actuators, \(u_{\mathrm{r},\ell}\) is the retained coordinate driven by the \(\ell\)th actuator, and \(\hat{u}_{\mathrm{r},\ell}\) is the displacement imposed on the specimen. Each actuator is assigned its own value of \(n_{\ell}\), which represents channels of different equivalent pure delay. A pure delay reproduces the phase lag of a channel rather than its bandwidth or its amplitude attenuation. The closed loop carrying these delays is shown in Fig. \ref{fig:rths-loop}. \rev{The numerical state supplies the target displacement of each actuated coordinate. The physical elastic and damping reactions return through the delayed channel terms. The additional regulator uses the numerical displacement and velocity and applies a delayed force at the same actuated coordinates.}

\begin{figure}[htbp]
	\centering
	% Selected option: Figure 1-D. Compare with TU thesis Fig. 2-1 (PDF p.21).
	\includegraphics[width=0.96\textwidth]{fig01_force_feedback.pdf}
	\caption{\rev{Virtual MDOF RTHS model with delayed physical reactions and an additional delayed force feedback.}}
	\label{fig:rths-loop}
\end{figure}

The retained coordinates of the physical substructure are the only coordinates that the transfer system imposes, and each of them is driven by one actuator. The condensed coordinates are neither commanded nor measured, and their response is reconstructed from the recovery relation of the condensation. The two condensation methods reconstruct this response in different ways.

The reconstructed response follows from Eqs. \eqref{eq:guyan-coordinate-relation} and \eqref{eq:cb-internal-recovery} as

\vspace{-3.5em}
\begin{subequations}
	\label{eq:delay-recovery}
	\begin{align}
		\bm{u}_{\mathrm{c}}^{\mathrm{rec}}(t)&=\bm{\Psi}_{\mathrm{G}}\hat{\bm{u}}_{\mathrm{r}}(t),
		\label{eq:guyan-delay-recovery}\\
		\bm{u}_{\mathrm{c}}^{\mathrm{rec}}(t)&=\bm{\Psi}_{\mathrm{G}}\hat{\bm{u}}_{\mathrm{r}}(t)
		+\bm{\Phi}_{d}\bm{\eta}(t),
		\label{eq:cb-delay-recovery}
	\end{align}
\end{subequations}
\vspace{-3.5em}

for Guyan and Craig--Bampton condensation, respectively. Every entry of \(\hat{\bm{u}}_{\mathrm{r}}\) is a delayed quantity. The Guyan reconstruction is therefore built from delayed quantities alone, and \(\bm{\Psi}_{\mathrm{G}}\) redistributes the delayed channels over all condensed coordinates. The fixed-interface modal coordinates \(\bm{\eta}\) are solved within the reduced model and do not pass through the transfer system. \rev{The modal contribution in Eq. \eqref{eq:cb-delay-recovery} is not directly delayed by an actuator; its evolution remains coupled to the delayed retained coordinates through the reduced equations.} \rev{The two methods therefore retain the same number of delay channels while representing their coupling to the internal response differently.}

The formulation is developed for a virtual RTHS, in which \(\bm{\eta}\) of the physical substructure is a state of the simulated specimen model.

\subsection{Discrete-time closed-loop pole criterion with LQR control}
\label{sec:pole-criterion}

The assembled reduced model of Eq. \eqref{eq:assembled-reduced} is taken as the starting point. The actuator delay affects only the terms acting on the coordinates imposed by the transfer system, and each channel carries its own delay. The assembled damping and stiffness matrices are therefore split channel by channel, and the equation of motion of the assembled reduced system is

\vspace{-3.5em}
\begin{equation}
	\mathbf{M}_{\mathrm{red}}\ddot{\bm{q}}(t)+\mathbf{C}_{0}\dot{\bm{q}}(t)+\mathbf{K}_{0}\bm{q}(t)
	+\sum_{\ell=1}^{n_{\mathrm{a}}}
	\left[
	\mathbf{C}_{\ell}\dot{\bm{q}}(t-\tau_{\ell})
	+\mathbf{K}_{\ell}\bm{q}(t-\tau_{\ell})
	\right]
	=\bm{f}_{\mathrm{red}}(t),
	\label{eq:delayed-eom}
\end{equation}
\vspace{-3.5em}

\noindent where \(\bm{q}(t)\) collects the \(n_{q}\) generalized coordinates of the assembled reduced model, \(\mathbf{C}_{\ell}\) and \(\mathbf{K}_{\ell}\) are of the same size as \(\mathbf{C}_{\mathrm{red}}\) and \(\mathbf{K}_{\mathrm{red}}\) and retain only the columns of the reduced physical substructure that act on the coordinate driven by the \(\ell\)th actuator, the remaining columns being zero, and \(\mathbf{C}_{0}=\mathbf{C}_{\mathrm{red}}-\sum_{\ell}\mathbf{C}_{\ell}\) and \(\mathbf{K}_{0}=\mathbf{K}_{\mathrm{red}}-\sum_{\ell}\mathbf{K}_{\ell}\) contain the remaining terms. This split makes the difference described in Section \ref{sec:multiple-delay} explicit. The reduced physical substructure of the Guyan model retains the actuated coordinates alone, and every one of its columns is therefore delayed. The Craig--Bampton model also carries fixed-interface modal coordinates driven by no actuator, and the columns associated with them remain in \(\mathbf{C}_{0}\) and \(\mathbf{K}_{0}\) without delay.

The force returned to the loop by the physical substructure is its elastic and damping reaction to the motion imposed by the transfer system. Each actuated channel therefore enters Eq. \eqref{eq:delayed-eom} with its own delay as a scalar factor, and the delays remain independent of each other. The inertia of the physical substructure is advanced numerically within the assembled state, and the assembled mass matrix carries no delay factor. A term \(\mathbf{M}_{\ell}\ddot{\bm{q}}(t-\tau_{\ell})\) would be added to Eq. \eqref{eq:delayed-eom} once the returned force contained the inertia produced by the delayed motion.

The assembled system of Eq. \eqref{eq:delayed-eom} is advanced with the CR integration algorithm \rev{\cite{ChenCR2008}}, and the sampling period equals the integration time step. A quantity written with the argument \(z\) denotes the \(z\) transform of the corresponding time-domain quantity in the following. The displacement and velocity recursions of the algorithm share one matrix coefficient. Eliminating the acceleration between them and taking the \(z\) transform relates the velocity and the acceleration to the displacement by

\vspace{-3.5em}
\begin{equation}
	\begin{gathered}
		\dot{\bm{q}}(z)=\frac{z-1}{z\Delta t}\bm{q}(z),\\
		\mathbf{M}_{\mathrm{red}}\ddot{\bm{q}}(z)=
		\left(\mathbf{M}_{\mathrm{red}}
		+\tfrac{1}{2}\Delta t\,\mathbf{C}_{\mathrm{red}}
		+\tfrac{1}{4}\Delta t^{2}\mathbf{K}_{\mathrm{red}}\right)
		\frac{(z-1)^{2}}{z\Delta t^{2}}\bm{q}(z),
	\end{gathered}
	\label{eq:z-derivatives}
\end{equation}
\vspace{-3.5em}

\noindent where the bracket is the mass term of the CR algorithm, formed from the assembled matrices before the channel split of Eq. \eqref{eq:delayed-eom}, since that split describes the signal path of the interface restoring force rather than the time integration. A delay of \(n_{\ell}\) steps becomes the scalar factor \(z^{-n_{\ell}}\). The delayed terms are taken from the stored response history, and the recursion therefore remains explicit in them.

A linear quadratic regulator acts on the actuated coordinates and \rev{supplies an additional generalized force at those coordinates through the delayed feedback path.} Its design model is the assembled reduced model statically condensed onto the actuated coordinates through Eqs. \eqref{eq:guyan-coordinate-relation} and \eqref{eq:guyan-transformation}, with the regulator force as its input. Minimizing a quadratic cost of the design-model states and of the control input, weighted by \(\mathbf{Q}\) and \(\mathbf{R}\), gives the control law

\vspace{-3.5em}
\begin{equation}
	\bm{f}_{\mathrm{L}}(t)=
	-\mathbf{K}_{x}\mathbf{S}_{\mathrm{a}}\bm{q}(t)
	-\mathbf{K}_{v}\mathbf{S}_{\mathrm{a}}\dot{\bm{q}}(t),
	\label{eq:lqr-law}
\end{equation}
\vspace{-3.5em}

\noindent where \(\mathbf{S}_{\mathrm{a}}\) is the Boolean matrix of size \(n_{\mathrm{a}}\times n_{q}\) that extracts the actuated coordinates from \(\bm{q}\), and \(\mathbf{K}_{x}\) and \(\mathbf{K}_{v}\) are the \(n_{\mathrm{a}}\times n_{\mathrm{a}}\) displacement and velocity parts of the gain.

The regulator output is a generalized force. \rev{Its applied value is \(\bm f_{\mathrm{L,applied}}(z)=\mathbf S_{\mathrm a}^{\mathsf T}\mathbf H_{\mathrm a}(z)\bm f_{\mathrm L}(z)\), where \(\mathbf H_{\mathrm a}(z)=\mathrm{diag}(z^{-n_1},\ldots,z^{-n_{\mathrm a}})\). This force enters the load summation of the numerical update, as shown in Fig. \ref{fig:rths-loop}; the gain matrices provide the corresponding feedback stiffness and damping.} The design model has \(2n_{\mathrm{a}}\) states against \(2n_{q}\) of the assembled model, and Eq. \eqref{eq:lqr-law} is accordingly an output feedback on the actuated coordinates rather than a full-state regulator of the assembled model. Each model supplies its own design model, and equal \(\mathbf{Q}\) and \(\mathbf{R}\) therefore give different gains. A stability domain computed in this way carries the regulator produced by the model itself.
Substituting Eqs. \eqref{eq:z-derivatives} and \eqref{eq:lqr-law} into Eq. \eqref{eq:delayed-eom} gives the closed-loop system in the \(z\) domain as \(\mathbf{Z}(z)\bm{q}(z)=\bm{f}_{\mathrm{red}}(z)\), with

\vspace{-3.5em}
\begin{equation}
	\begin{aligned}
		\mathbf{Z}(z)={}&
		\frac{(z-1)^{2}}{z\Delta t^{2}}
		\left(\mathbf{M}_{\mathrm{red}}
		+\tfrac{1}{2}\Delta t\,\mathbf{C}_{\mathrm{red}}
		+\tfrac{1}{4}\Delta t^{2}\mathbf{K}_{\mathrm{red}}\right)
		+\frac{z-1}{z\Delta t}\mathbf{C}_{0}+\mathbf{K}_{0}\\
		&+\sum_{\ell=1}^{n_{\mathrm{a}}}z^{-n_{\ell}}
		\left[
		\frac{z-1}{z\Delta t}\mathbf{C}_{\ell}+\mathbf{K}_{\ell}
		\right]
		+\mathbf{S}_{\mathrm{a}}^{\mathsf{T}}\mathbf{H}_{\mathrm{a}}(z)
		\left[
		\mathbf{K}_{x}+\frac{z-1}{z\Delta t}\mathbf{K}_{v}
		\right]\mathbf{S}_{\mathrm{a}},
	\end{aligned}
	\label{eq:z-matrix}
\end{equation}
\vspace{-3.5em}

\noindent where \(\mathbf{H}_{\mathrm{a}}(z)\) is the diagonal matrix whose \(\ell\)th entry is the delay factor \(z^{-n_{\ell}}\) of the \(\ell\)th channel. The poles of the closed-loop system are the roots of

\vspace{-3.5em}
\begin{equation}
	\det\mathbf{Z}(z)=0.
	\label{eq:characteristic-equation}
\end{equation}
\vspace{-3.5em}

The negative powers of \(z\) in Eq. \eqref{eq:z-matrix} are cleared by multiplying Eq. \eqref{eq:characteristic-equation} by a sufficient power of \(z\), and the roots introduced at the origin by this operation are discarded.

The continuous and the discrete planes are related by \(z=\mathrm{e}^{s\Delta t}\), and the left half of the \(s\) plane therefore maps onto the interior of the unit circle of the \(z\) plane. The largest modulus among the roots of Eq. \eqref{eq:characteristic-equation} is written as \(\rho_{\mathrm{cl}}\). The system is stable when \(\rho_{\mathrm{cl}}<1\), unstable when \(\rho_{\mathrm{cl}}>1\), and on the stability boundary when \(\rho_{\mathrm{cl}}=1\).

For a given condensed model, \(\rho_{\mathrm{cl}}\) depends on the delay of each actuator through the factors \(z^{-n_{\ell}}\). The stability domain is obtained by evaluating \(\rho_{\mathrm{cl}}\) over combinations of the integer delays \(n_{\ell}\) and by locating the contour on which \(\rho_{\mathrm{cl}}\) equals unity. The delays are expressed in multiples of the sampling period, and the resulting domain gives the combinations of actuator delays for which the RTHS system remains stable.




\section{Benchmark structure and analysis setup} \label{sec4}

The formulations of Sections \ref{sec2} and \ref{sec3} are applied to a benchmark frame whose physical substructure contains more interface coordinates than the transfer system can impose independently. The study is a virtual RTHS, in which both substructures are simulated and the two-channel actuation limit is imposed as a modelling constraint. The accuracy and the stability of a condensed model depend on the extent of the physical substructure and on the choice of the actuated coordinates, and two substructure divisions with different condensation ratios are therefore examined. In this section, the benchmark structure and its finite element idealization are described first. The two substructure divisions are then defined together with the retained and condensed coordinates of each substructure. The excitations, the controller and delay settings, and the indices adopted for the modelling accuracy and the stability are given last.

\subsection{Benchmark structure and finite element idealization}
\label{sec:benchmark}

The reference structure is a simplified planar model built from the geometry and \rev{section arrangement} of the multi-axial RTHS benchmark frame of Condori Uribe et al. \cite{Condori2023}. The three bays are 762 mm wide and the three storeys are 635 mm high. The columns are W5\(\times\)16 sections of A992 Gr.50 steel and the beams are a custom I section of A36 steel. The geometry and the section dimensions are shown in Fig. \ref{fig:benchmark-geometry}, and the material properties are listed in Table \ref{tab:materials}. \rev{The numerical model uses \(E=206\) GPa for both member types and the section properties stated below.}

\begin{figure}[htbp]
	\centering
	% Selected option: Figure 3-C. Compare with TU thesis Fig. 2-3 (PDF p.28).
	\includegraphics[width=0.89\textwidth]{fig02_geometry.pdf}
	\caption{\rev{Geometry and member sections of the reference frame; dimensions in mm. The numerical section properties are specified in the text.}}
	\label{fig:benchmark-geometry}
\end{figure}

\begin{table}[htbp]
	\centering
	\caption{Material properties of the benchmark frame.}
	\label{tab:materials}
	\begin{tabular}{lccccc}
		\toprule
		Steel grade & \(f_{\mathrm{y}}\) (MPa) & \(f_{\mathrm{u}}\) (MPa) & \(E\) (GPa) & \(\nu\) & \(\rho\) (kg/m\(^{3}\))\\
		\midrule
		A36 (beams)      & 250 & 400--550 & \rev{206} & 0.3 & 7850\\
		A992 Gr.50 (columns) & 345 & 450--550 & \rev{206} & 0.3 & 7850\\
		\bottomrule
	\end{tabular}
\end{table}

In the finite element idealization, each node of the frame carries two in-plane translations and one rotation about the out-of-plane axis. The frame responds mainly in the horizontal direction under a base excitation. The axial deformation of the members is therefore neglected, the vertical displacement is taken as zero, and the horizontal displacement is assumed identical for all nodes of a storey. The idealized model of Fig. \ref{fig:dof-idealization} has fifteen DOFs, which are one horizontal displacement per storey and one rotation at each of the four nodes of each storey. Among them, \(\psi_{1}\), \(\psi_{6}\) and \(\psi_{11}\) are the horizontal displacements of the first, second and third storey, and the remaining coordinates are the nodal rotations of the corresponding storey.

\rev{The numerical frame fixes all four column bases. Its effective section properties are \(A_{\mathrm c}=1.0774172\times10^{-3}\,\mathrm{m}^2\), \(I_{\mathrm c}=1.0489032\times10^{-6}\,\mathrm{m}^4\), \(A_{\mathrm b}=6.1096652\times10^{-4}\,\mathrm{m}^2\), and \(I_{\mathrm b}=2.5523311\times10^{-7}\,\mathrm{m}^4\). A diagonal lumped mass matrix assigns 6164.831 kg to each storey translation and rotational inertias of 51.088, 67.887, 67.887 and 51.088 kg\,m\(^{2}\) to its four nodes. These masses use an effective density multiplier of 190 relative to the steel density in Table \ref{tab:materials}; the multiplier sets the numerical mass level. The matrices and coordinate order are supplied with the paper.}

\begin{figure}[htbp]
	\centering
	% Selected option: Figure 4-D.
	% Panel (a): TU thesis Fig. 2-4 (PDF p.29); panel (b): Fig. 2-5 (PDF p.30).
	\includegraphics[width=0.77\textwidth]{fig03_dofs.pdf}
	\caption{\rev{Degrees of freedom and node numbering of the idealized numerical model.}}
	\label{fig:dof-idealization}
\end{figure}

The full-order equation of motion is Eq. \eqref{eq:full-order-eom}. The damping matrix is proportional to the mass and stiffness matrices, \(\mathbf{C}=a_{0}\mathbf{M}+a_{1}\mathbf{K}\), and the coefficients \(a_{0}\) and \(a_{1}\) are obtained from a damping ratio of 5\% at the first two modes. \rev{The coefficients are \(a_0=1.3184625\,\mathrm{s}^{-1}\) and \(a_1=1.3223424\times10^{-3}\,\mathrm{s}\); the reduced damping matrices are obtained from this same damping matrix. The first five natural frequencies of the idealized model are 2.7074 Hz, 9.3284 Hz, 18.4427 Hz, 22.6736 Hz and 24.2277 Hz. The first five modes account for 98.06\%} of the effective modal mass and therefore characterize the dominant seismic response of the structure.

\subsection{Substructure division and selection of retained DOFs}
\label{sec:division}

The position of the boundary between the physical and the numerical substructures affects both the dynamic response and the accuracy of the condensation. Two divisions are compared and both are shown in Fig. \ref{fig:division}. In each of them the outer bay of the frame is implemented as the physical substructure and the remaining part forms the numerical substructure. The outer bay is selected because a division at an inner bay would split the numerical substructure into two separate parts.

In division I the physical substructure covers the first two storeys of the outer bay together with the connected beams and columns. The lower storeys carry the largest storey shear of the first two modes, and the interface is close to the actuator locations, which shortens the loading path. The physical substructure holds about 40\% of the DOFs of the whole structure.

In division II the physical substructure covers all three storeys of the outer bay. More DOFs are kept at the interface, which represents the modal coupling between the storeys more completely, and the physical substructure holds about 60\% of the DOFs of the whole structure.

Three requirements govern the selection of the retained coordinates. The coordinates carrying the dominant response remain explicit, the condensed response is recoverable through the relations of Section \ref{sec2}, and the reduction is large enough to be of practical use. The horizontal storey displacements govern the global response of a frame under a base excitation and carry most of the energy of the low-order modes. The nodal rotations mainly represent local bending and are coupled to the horizontal displacements through the off-diagonal terms of the stiffness matrix, and they are condensed accordingly. Condensing all rotations would shift the higher modes, and several rotations of the numerical substructure are therefore kept.

Both divisions are limited to two independent actuation channels. Division I directly actuates both horizontal coordinates of its physical substructure. Division II has three horizontal coordinates under the same two-channel constraint, and the middle-storey coordinate \(\psi_{6}\) is therefore condensed and reconstructed from the retained coordinates. Division II is the more demanding case, since a coordinate carrying a substantial share of the storey response is no longer imposed directly. The retained horizontal coordinates of the physical substructure are the coordinates imposed by the two actuators in both divisions. The retained set and the actuated set thus coincide, \rev{and these channel locations are used to select the coordinates monitored in the accuracy comparison and delayed in the stability formulation.}

Three rotations of the numerical substructure are retained as internal master coordinates, which limits the shift of its higher modes. They leave the interface coordinates condensed by the physical substructure outside the coupling set, where the compatibility is not restored. The retained and condensed coordinates of both substructures are listed in Table \ref{tab:dof-sets} with the coordinate indices of Fig. \ref{fig:dof-idealization}, and a coordinate lying on the interface appears in both substructures.

\begin{figure}[htbp]
\centering
\includegraphics[width=.98\textwidth,height=.74\textheight,keepaspectratio]{fig04_divisions.pdf}
\caption{\rev{Substructure divisions and local retained coordinates. (a) Division I. (b) Division II. The physical and numerical coordinate sets are listed in Table \ref{tab:dof-sets}.}}
\label{fig:division}
\end{figure}

\begin{table}[htbp]
	\centering
	\caption{Retained and condensed coordinates of the two substructure divisions.}
	\label{tab:dof-sets}
	\begin{tabular}{llll}
		\toprule
		Division & Substructure & Retained coordinates & Condensed coordinates\\
		\midrule
		I  & Numerical & \(\psi_{1},\psi_{4},\psi_{6},\psi_{9},\psi_{11},\psi_{14}\)
		& \(\psi_{3},\psi_{5},\psi_{7},\psi_{8},\psi_{10},\psi_{12},\psi_{13},\psi_{15}\)\\
		I  & Physical  & \(\psi_{1},\psi_{6}\)
		& \(\psi_{2},\psi_{3},\psi_{7},\psi_{8}\)\\
		\midrule
		II & Numerical & \(\psi_{1},\psi_{4},\psi_{6},\psi_{9},\psi_{11},\psi_{14}\)
		& \(\psi_{3},\psi_{5},\psi_{8},\psi_{10},\psi_{13},\psi_{15}\)\\
		II & Physical  & \(\psi_{1},\psi_{11}\)
		& \(\psi_{2},\psi_{3},\psi_{6},\psi_{7},\psi_{8},\psi_{12},\psi_{13}\)\\
		\bottomrule
	\end{tabular}
\end{table}

The three fixed-interface modes of lowest frequency are retained for each substructure in the Craig--Bampton models, which gives the two divisions the same Craig--Bampton order. Assembly through the two interface coordinates of Table \ref{tab:dof-sets} gives a Guyan model with 6 generalized coordinates and a Craig--Bampton model with 12 in both divisions. \rev{These substructure-level configurations have the same reduced order in the two divisions.}

\rev{For the delay-free accuracy comparisons, the assembled fifteen-DOF frame is reduced directly. The retained physical coordinates are \(\{\psi_1,\psi_6,\psi_{11},\psi_4,\psi_9,\psi_{14}\}\) in division I and \(\{\psi_1,\psi_{11},\psi_4,\psi_9,\psi_{14}\}\) in division II. Three fixed-interface modes enrich each full-frame Craig--Bampton basis, giving orders 6/9 for Guyan/Craig--Bampton in division I and 5/8 in division II. Each response is expanded uniquely to the fifteen physical coordinates. These models evaluate the effect of the retained set on structural response; the local substructure coupling sets remain those in Table \ref{tab:dof-sets}.}

\rev{The Guyan accuracy implementation substitutes the static recovery relation into the retained equations, using \(\mathbf M_{\mathrm G}=\mathbf M_{\mathrm{rr}}+\mathbf M_{\mathrm{rc}}\bm\Psi_{\mathrm G}\), with analogous expressions for \(\mathbf C_{\mathrm G}\) and \(\mathbf K_{\mathrm G}\), and the load \(\mathbf T_{\mathrm G}^{\mathsf T}\bm f\). The Craig--Bampton matrices use Eq. \eqref{eq:general-projection}. Thus the accuracy comparison also includes this distinction in the Guyan matrix construction.}

\subsection{Simulation setup and evaluation indices}
\label{sec:setup}

\rev{The full-order model and the two full-frame condensed models are implemented as three parallel state-space systems in Simulink. The accuracy study uses zero actuator delay and zero regulator force, \(\bm f_{\mathrm L}=\bm0\), with no measurement noise. The equivalent seismic load is formed from the full mass matrix and projected to the reduced coordinates. A fixed-step fourth-order Runge--Kutta solver advances the systems at \(\Delta t=1/1024\) s over 40 s, and the transformations recover the three storey displacements. Differences from the full-order response therefore measure the accuracy of the stated reduction implementations.}

Two excitations are adopted. The first is the El Centro record scaled by a factor of 0.40\rev{, using the supplied north--south acceleration series in m/s\(^{2}\), with its 0.02 s samples linearly interpolated during integration. The response is evaluated with the linear model of Eq. \eqref{eq:full-order-eom}.} The second is a chirp signal whose frequency increases linearly from 0.1 Hz to 10 Hz over 40 s. \rev{Its acceleration is \(a_{\mathrm g}(t)=\sin[2\pi(0.1)t+\pi(9.9/40)t^2]\,\mathrm{m/s^2}\), with unit amplitude and zero initial phase. Both input histories are shown in Fig. \ref{fig:inputs}.} The sweep covers the first two natural frequencies of the benchmark and follows the accuracy of the condensed model as the excitation frequency approaches and then exceeds the fundamental frequency.

\begin{figure}[htbp]
\centering
\includegraphics[width=.98\textwidth]{fig11_excitation.pdf}
\caption{\rev{Ground-acceleration inputs. (a) Scaled El Centro record. (b) Linear chirp.}}
\label{fig:inputs}
\end{figure}

The stability analysis adopts the closed loop of Section \ref{sec3}. The sampling period is \(\Delta t=1/1024\) s and equals the integration time step. Two actuators drive the retained horizontal coordinates of the physical substructure. Actuator 1 drives \(\psi_{1}\) with a delay \(\tau_{1}\) in both divisions, while actuator 2 drives \(\psi_{6}\) in division I and \(\psi_{11}\) in division II, both with a delay \(\tau_{2}\). The two delays are treated as independent, which represents channels of different equivalent pure delay. The weighting matrices of the regulator are \(\mathbf{Q}=\mathrm{diag}(10^{6},10^{6},10^{4},10^{4})\) and \(\mathbf{R}=\mathrm{diag}(10^{-2},10^{-2})\), where the first two entries of \(\mathbf{Q}\) weight the displacements of the two actuated coordinates and the last two weight their velocities.

Three indices are adopted for the accuracy of the condensed models. The first is the relative error of the natural frequencies, \(e_{f,i}=|f_{i}^{\mathrm{full}}-f_{i}^{\mathrm{red}}|/f_{i}^{\mathrm{full}}\times100\%\), where \(f_{i}^{\mathrm{full}}\) and \(f_{i}^{\mathrm{red}}\) are the \(i\)th natural frequency of the full-order and of the condensed model. The second is the modal assurance criterion,

\vspace{-3.5em}
\begin{equation}
	\mathrm{MAC}_{i}=
	\frac{\left[\left(\bm{\varphi}_{i}^{\mathrm{full}}\right)^{\mathsf{T}}
		\bm{\varphi}_{i}^{\mathrm{red}}\right]^{2}}
	{\left\|\bm{\varphi}_{i}^{\mathrm{full}}\right\|^{2}
		\left\|\bm{\varphi}_{i}^{\mathrm{red}}\right\|^{2}},
	\label{eq:mac}
\end{equation}
\vspace{-3.5em}

\noindent where \(\bm{\varphi}_{i}^{\mathrm{red}}\) is the \(i\)th mode shape of the condensed model and \(\bm{\varphi}_{i}^{\mathrm{full}}\) is the \(i\)th mode shape of the full-order model restricted to the retained coordinates, which places the two vectors on the same coordinate set. The retained-coordinate part of the reduced mode shape is taken for the Craig--Bampton models. \rev{Translations are expressed in metres and rotations in radians, with no additional coordinate scaling in this retained-coordinate MAC. A value close to unity indicates agreement under this coordinate convention.} The third is the normalized root mean square error of the time histories,

\vspace{-3.5em}
\begin{equation}
	\mathrm{NRMSE}=
	\frac{\displaystyle\sqrt{\frac{1}{T_{\mathrm{d}}}\int_{0}^{T_{\mathrm{d}}}
			\left[x^{\mathrm{full}}(t)-x^{\mathrm{red}}(t)\right]^{2}\mathrm{d}t}}
	{\max\left(x^{\mathrm{full}}\right)-\min\left(x^{\mathrm{full}}\right)}
	\times100\%,
	\label{eq:nrmse}
\end{equation}
\vspace{-3.5em}

\noindent where \(x^{\mathrm{full}}\) and \(x^{\mathrm{red}}\) are the horizontal storey displacements of the full-order and of the condensed model, and \(T_{\mathrm{d}}\) is the duration of the response. For a frequency band, the numerator is evaluated over the corresponding time interval and divided by the duration of that interval, which places bands of different length on the same root mean square basis. The denominator is the peak-to-peak value of the full-order response over the whole record and is common to every band, and a band of small response amplitude therefore returns a small NRMSE.

For the chirp excitation, the NRMSE is evaluated separately over five frequency bands, \rev{with endpoints of 0.1, 1.9, 3.5, 5.4, 8.1 and 10 Hz. These fixed endpoints approximate multiples of the fundamental frequency, \(f_1=2.7074\) Hz, and are used exactly in the error integration.} The instantaneous frequency of the sweep is \(f(t)=0.1+0.2475t\). The band containing the fundamental frequency therefore spans the interval from about 7.3 s to 13.7 s, and the sweep crosses \(f_{1}\) at about 10.5 s. The bands are centred on the fundamental frequency, where the condensation error is most sensitive to the excitation frequency. The response within a band carries the decay of the preceding band, and \rev{the band values accordingly track the variation of error across the sweep.}

\rev{The change in low-order modal distribution is described by a modal redistribution ratio, which combines the relative changes at the actuated coordinates.} With the mode shapes of the full-order model normalized with respect to the mass matrix, the modal share of coordinate \(j\) is expressed as

\vspace{-3.5em}
\begin{equation}
	\begin{gathered}
		\gamma_{i}=
		\frac{\bm{\varphi}_{i}^{\mathsf{T}}\mathbf{M}\bm{\Gamma}}
		{\bm{\varphi}_{i}^{\mathsf{T}}\mathbf{M}\bm{\varphi}_{i}},
		\qquad
		p_{i}=\frac{\gamma_{i}^{2}}{\displaystyle\sum_{k=1}^{m}\gamma_{k}^{2}},\\[2pt]
		E_{j}=\sum_{i=1}^{m}p_{i}\,
		\frac{m_{j}\varphi_{j,i}^{2}}{\displaystyle\sum_{l=1}^{n}m_{l}\varphi_{l,i}^{2}},
	\end{gathered}
	\label{eq:dof-energy}
\end{equation}
\vspace{-3.5em}

\noindent where \(\gamma_{i}\) is the modal participation factor of the \(i\)th mode, \(p_{i}\) is the corresponding modal weight, \(E_{j}\) is the modal share of coordinate \(j\), \(\bm{\varphi}_{i}\) is the \(i\)th mode shape and \(\varphi_{j,i}\) is its \(j\)th entry, \(\bm{\Gamma}\) is the influence vector of Eq. \eqref{eq:full-order-eom}, \(m_{j}\) is the diagonal entry of the mass matrix associated with coordinate \(j\), \(n\) is the number of DOFs of the full-order model, and \(m\) is the number of modes considered, taken as \(m=5\) for the frame of Section \ref{sec:benchmark}. The mode shapes of a condensed model are expanded to the full coordinate set through its own transformation matrix before Eq. \eqref{eq:dof-energy} is applied, \rev{and mass-normalized again using the full mass matrix. The participation weights are recomputed for each model, so both shares are evaluated on the same fifteen physical coordinates and mass convention.} The modal redistribution ratio is then obtained as

\vspace{-3.5em}
\begin{equation}
	\Delta E=\sum_{k=1}^{n_{\mathrm{a}}}
	\frac{E^{\mathrm{red}}(d_{k})-E^{\mathrm{full}}(d_{k})}{E^{\mathrm{full}}(d_{k})}.
	\label{eq:energy-change}
\end{equation}
\vspace{-3.5em}

\noindent where \(d_{k}\) is the \(k\)th actuated coordinate of the physical substructure, and \(E^{\mathrm{full}}\) and \(E^{\mathrm{red}}\) are the modal shares before and after condensation. Equation \eqref{eq:energy-change} sums the relative change of the modal share of each actuated coordinate. \rev{The terms have separate denominators and retain their signs; their sum is a modal-share indicator rather than a total energy-transfer measure or a stability criterion.} The index is built from the diagonal entries of the mass matrix and is applied to comparisons on the same physical coordinate set.




\section{Results and discussion} \label{sec5}

The two condensed models of each division are compared with the full-order reference model of Section \ref{sec:benchmark}, following the order in which condensation acts on an RTHS. The modal properties of the reduced models are examined first, since a shift of the natural frequencies changes the response at every excitation level. The response accuracy under the two excitations is evaluated next \rev{for the full-frame reduction implementations defined in Section \ref{sec:division}.} \rev{The substructure-level delay domains are then compared for pairs of actuator delays.} \rev{The modal redistribution ratio finally describes how the low-order modal shares at the actuated coordinates change under condensation.}

\subsection{Modal accuracy}
\label{sec:modal-accuracy}

The relative errors of the first two natural frequencies and the corresponding MAC values are listed in Table \ref{tab:modal}. The chirp excitation reaches 10 Hz, and the first two modes at \rev{2.7074 Hz and 9.3284 Hz} are therefore the modes within the excitation band, to which the comparison is restricted. Craig--Bampton condensation reproduces the natural frequencies of the full-order model in both divisions, with a largest error of \rev{0.284\%}. Guyan condensation is accurate for the first mode of division I, where the error is 0.648\%, while the error of the second mode reaches 5.411\%. Both modes are affected in division II, with errors of \rev{19.040\% and 30.574\%}.

\begin{table}[htbp]
	\centering
	\caption{Relative error of the natural frequencies and MAC values of the condensed models.}
	\label{tab:modal}
	\begin{tabular}{llcccc}
		\toprule
		& & \multicolumn{2}{c}{Frequency error (\%)} & \multicolumn{2}{c}{MAC}\\
		\cmidrule(lr){3-4}\cmidrule(lr){5-6}
		Division & Mode & Guyan & Craig--Bampton & Guyan & Craig--Bampton\\
		\midrule
		I  & 1 & 0.648  & 0.003 & 1.0000 & 1.0000\\
		I  & 2 & 5.411  & 0.284 & 0.9998 & 1.0000\\
		\midrule
		II & 1 & \rev{19.040} & \rev{0.024} & 0.9962 & 1.0000\\
		II & 2 & \rev{30.574} & \rev{0.108} & 0.9255 & \rev{1.0000}\\
		\bottomrule
	\end{tabular}
\end{table}

The two divisions differ in that division II condenses a horizontal storey coordinate \rev{from the full-frame accuracy model. Once the middle-storey coordinate leaves the retained set, the static recovery is controlled by the retained lower- and upper-storey translations and the selected rotations. In the retained-equation Guyan implementation, the diagonal full mass matrix gives \(\mathbf M_{\mathrm G}=\mathbf M_{\mathrm{rr}}\), while the stiffness and load retain the static recovery contribution. The resulting change in the balance of stiffness and inertia raises the frequencies in this case. Craig--Bampton condensation includes internal modal inertia through its fixed-interface modes and retains the low-order frequencies much more closely. Figure \ref{fig:horizontal-modes} displays the actual frequencies together with the restored horizontal mode shapes.}

\begin{figure}[htbp]
\centering
\includegraphics[width=.98\textwidth]{fig12_horizontal_modes.pdf}
\caption{\rev{First two horizontal mode shapes expanded to the three storeys. Each shape is normalised by its largest horizontal component; the corresponding natural frequencies (Hz) are listed beneath each panel.}}
\label{fig:horizontal-modes}
\end{figure}

The MAC values behave differently. All four models keep the MAC above 0.92, and the lowest value of 0.9255 belongs to the Guyan model of division II, whose second natural frequency is in error by \rev{30.574\%}. The mode shape components on the retained coordinates are therefore preserved where the frequencies are not. \rev{The static deformation pattern can resemble the retained components of a mode while the reduced mass representation changes its frequency. The restored second mode in Fig. \ref{fig:horizontal-modes} also shows the larger spatial deviation of the division II Guyan model.} A MAC evaluated on the retained coordinates alone is accordingly insufficient to accept a condensed model, and the frequency error is the modal quantity to be monitored alongside it.

\FloatBarrier
\subsection{Response accuracy under earthquake and chirp excitation}
\label{sec:response-accuracy}

Figs. \ref{fig:eq-div1} and \ref{fig:eq-div2} compare the horizontal storey displacements of the two condensed models with those of the full-order model under the El Centro record. The first and the third storey are shown. The first storey is actuated in both divisions and allows the same coordinate to be followed throughout, and the third storey carries the largest lateral response of the low-order modes. \rev{The second-storey response is examined separately in Fig. \ref{fig:middle-eq}, because its reconstruction is the critical additional approximation in division II.} In division I both condensed models follow the reference over the whole record, and the deviation of the Guyan model appears only in the enlarged window between 10 s and 11 s. In division II the Guyan response departs from the reference over the same window in both amplitude and shape, whereas the Craig--Bampton response stays close to it. The window between 21.5 s and 22.5 s carries a weaker excitation and is reproduced by both models in both divisions, and the error of the Guyan model is therefore not uniform in time.

\begin{figure}[htbp]
	\centering
	% Calculation-result plot: source/results mapping in figure/results_v2.
	% Compare with TU thesis Fig. 3-6 (PDF p.42) and Fig. 3-8 (PDF p.43).
	\includegraphics[width=0.98\textwidth]{fig06_eq_div1.pdf}
	\caption{Horizontal storey displacement of division I under the El Centro record. (a) First storey. (b) Third storey.}
	\label{fig:eq-div1}
\end{figure}

\begin{figure}[htbp]
	\centering
	% Calculation-result plot: source/results mapping in figure/results_v2.
	% Compare with TU thesis Fig. 3-9 (PDF p.44) and Fig. 3-10 (PDF p.45).
	\includegraphics[width=0.98\textwidth]{fig07_eq_div2.pdf}
	\caption{Horizontal storey displacement of division II under the El Centro record. (a) First storey. (b) Third storey.}
	\label{fig:eq-div2}
\end{figure}

\rev{Figure \ref{fig:middle-eq} evaluates the middle-storey displacement recovered from the division II full-frame transformations. The Guyan NRMSE is 10.798\%, whereas the Craig--Bampton value is 0.031\% over the 40 s record. The separate error traces show that the small Craig--Bampton error persists through the strong-response interval, rather than following only the overall amplitude envelope. This comparison directly tests recovery of the omitted horizontal coordinate.}

\begin{figure}[htbp]
\centering
\includegraphics[width=.98\textwidth]{fig13_middle_eq.pdf}
\caption{\rev{Reconstructed middle-storey displacement in division II under El Centro excitation. (a) Response comparison. (b) Guyan error. (c) Craig--Bampton error. Errors are measured relative to the full-order response; the two error panels use different vertical scales.}}
\label{fig:middle-eq}
\end{figure}

Figs. \ref{fig:chirp-div1} and \ref{fig:chirp-div2} show the response under the chirp excitation, in which the instantaneous frequency rises linearly from 0.1 Hz to 10 Hz over the 40 s of the record. The sweep crosses the fundamental frequency at about 10.5 s and reaches three times the fundamental frequency at about 32 s. In the low-frequency part of the sweep both condensed models reproduce the reference in both divisions.

In division I the Guyan response deviates from the reference around the \rev{full-order first-storey response peak at 11.69 s}, and the deviation remains small. The resonant response builds up over several cycles, and this peak therefore lags the crossing of the fundamental frequency at about 10.5 s. The deviation grows again towards the end of the sweep near 38 s, where the instantaneous frequency is about 9.5 Hz and the excitation passes the second mode of the structure. The second natural frequency of this model is in error by 5.411\%, which places the second resonance at the wrong frequency and turns the deviation into an amplitude and phase error over that window. In division II the Guyan response is inaccurate over the whole resonance band, and its peaks are shifted in time against those of the reference, which is the time-domain signature of the frequency shift of Section \ref{sec:modal-accuracy}. Near the end of the sweep, the Guyan response of division II \rev{has peak-to-peak amplitudes of 46.1\% and 80.5\% of the reference at the first and third storeys, respectively, over 38.0--38.3 s.} The Craig--Bampton response follows the reference over the whole sweep in both divisions.

\begin{figure}[htbp]
	\centering
	% Calculation-result plot: source/results mapping in figure/results_v2.
	% Compare with unique thesis Fig. 3-11 (PDF p.46) and Fig. 3-13 (p.47--48).
	\includegraphics[width=0.98\textwidth]{fig08_chirp_div1.pdf}
	\caption{Horizontal storey displacement of division I under the chirp excitation. (a) First storey. (b) Third storey.}
	\label{fig:chirp-div1}
\end{figure}

\begin{figure}[htbp]
	\centering
	% Calculation-result plot: source/results mapping in figure/results_v2.
	% Compare with unique thesis Fig. 3-14 (PDF p.48 lower) and Fig. 3-15 (PDF p.49).
	\includegraphics[width=0.98\textwidth]{fig09_chirp_div2.pdf}
	\caption{Horizontal storey displacement of division II under the chirp excitation. (a) First storey. (b) Third storey.}
	\label{fig:chirp-div2}
\end{figure}

\rev{The corresponding middle-storey sweep response is shown in Fig. \ref{fig:middle-chirp}. The Guyan and Craig--Bampton full-record NRMSE values are 11.749\% and 0.034\%, respectively. The Guyan error concentrates around the resonance passages, consistent with its shifted natural frequencies. The Craig--Bampton reconstruction follows the middle-storey oscillations through both passages, showing that its improvement extends to the horizontal coordinate excluded from the retained set.}

\begin{figure}[htbp]
\centering
\includegraphics[width=.98\textwidth]{fig14_middle_chirp.pdf}
\caption{\rev{Reconstructed middle-storey displacement in division II under chirp excitation. (a) Response comparison. (b) Guyan error. (c) Craig--Bampton error. The two error panels use different vertical scales.}}
\label{fig:middle-chirp}
\end{figure}

Table \ref{tab:nrmse} gives the NRMSE of the first-storey horizontal displacement. Under the El Centro record the Guyan error is \rev{0.802\%} in division I and \rev{10.911\%} in division II, whereas the Craig--Bampton error stays below 0.03\% in both. The band decomposition of the chirp response locates the origin of the Guyan error. Below the fundamental frequency the Guyan error is \rev{0.027\% in division I and 0.719\% in division II}. It rises to \rev{2.010\% and 26.858\%} in the band containing the fundamental frequency, the latter being the largest value of the table. Above the fundamental frequency the division I error falls back to \rev{0.171\%} in the band from 5.4 Hz to 8.1 Hz and rises again to \rev{2.041\%} in the highest band, whose small absolute response places the corresponding NRMSE below the relative deviation of Fig. \ref{fig:chirp-div1} near 38 s. \rev{In division II the Guyan error drops to 0.129\% in the 5.4--8.1 Hz band and rises to 3.020\% in the final band. Craig--Bampton gives a maximum band NRMSE of 0.138\% in division I and 0.074\% in division II.}

\begin{table}[htbp]
	\centering
	\caption{NRMSE of the first-storey horizontal displacement (\%).}
	\label{tab:nrmse}
	\begin{tabular}{lcccc}
		\toprule
		& \multicolumn{2}{c}{Division I} & \multicolumn{2}{c}{Division II}\\
		\cmidrule(lr){2-3}\cmidrule(lr){4-5}
		Excitation & Guyan & Craig--Bampton & Guyan & Craig--Bampton\\
		\midrule
		El Centro           & \rev{0.802} & \rev{0.014}    & \rev{10.911} & \rev{0.029}\\
		\midrule
		Chirp, 0.1--1.9 Hz  & \rev{0.027} & $<$0.001 & \rev{0.719}  & \rev{0.001}\\
		Chirp, 1.9--3.5 Hz  & \rev{2.010} & 0.008    & \rev{26.858} & \rev{0.074}\\
		Chirp, 3.5--5.4 Hz  & \rev{0.631} & 0.003    & \rev{9.898} & \rev{0.023}\\
		Chirp, 5.4--8.1 Hz  & \rev{0.171} & \rev{0.007}    & \rev{0.129}  & \rev{0.004}\\
		Chirp, 8.1--10 Hz   & \rev{2.041} & \rev{0.138}    & \rev{3.020}  & \rev{0.053}\\
		\bottomrule
	\end{tabular}
\end{table}

Two observations follow. The accuracy of a condensed model depends on the excitation frequency, and the band around the fundamental frequency is where the two methods differ most. \rev{The Guyan errors differ by more than an order of magnitude between the divisions in the fundamental-frequency band. The accuracy models have six and five retained coordinates, respectively, and the loss of the middle-storey translation is the central change in division II. Figures \ref{fig:middle-eq} and \ref{fig:middle-chirp} show its effect directly at the reconstructed coordinate.}

\rev{Peak interstorey drifts provide a complementary measure of the recovered deformation pattern (Fig. \ref{fig:peak-drift}). Each drift is the relative displacement of adjacent floors divided by the 635 mm storey height, with zero base displacement. In division II, the Guyan third-storey peak drift rises from the reference 0.687\% to 0.798\% under El Centro excitation and from 1.784\% to 2.259\% under the chirp. The corresponding Craig--Bampton values are 0.689\% and 1.789\%. Thus the improved coordinate recovery also preserves the distribution of interstorey deformation, whereas an individual floor peak alone can conceal an error in the relative motion.}

\begin{figure}[htbp]
\centering
\includegraphics[width=.98\textwidth]{fig15_peak_drift.pdf}
\caption{\rev{Peak absolute interstorey drift profiles. (a,b) El Centro excitation. (c,d) Chirp excitation. Panels (a,c) correspond to division I, and panels (b,d) to division II.}}
\label{fig:peak-drift}
\end{figure}

\FloatBarrier
\subsection{Stability domains under multiple actuator delays}
\label{sec:stability-domains}

\rev{The reference delay-domain comparison is shown in Fig. \ref{fig:stability-domain}. The two axes give the channel delays in milliseconds, using the exact conversion \(\tau_\ell=1000n_\ell/1024\) ms from the integer delay steps. Each supplied boundary gives the largest stable second-channel delay for a prescribed first-channel delay, with the stable region on the side containing the origin. The Original boundary represents the unreduced fifteen-DOF reference with ideal coupling at the interface coordinates outside the two delayed channels. It provides the ideal-interface reference for this comparison.} \rev{These boundaries are retained from the original boundary dataset; the available full-frame accuracy matrices do not establish the closed-loop matrices and regulator gains behind these curves.}

\begin{figure}[htbp]
	\centering
	% Historical stability-boundary snapshot, plotting-level reproduction only.
	% Compare with TU thesis Fig. 4-4 and Fig. 4-5 (both on PDF p.70).
	\includegraphics[width=0.98\textwidth]{fig10_stability_domain.pdf}
	\caption{Stability domains in the plane of the two actuator delays. (a) Division I. (b) Division II.}
	\label{fig:stability-domain}
\end{figure}

Both condensations reduce the stability domain relative to the unreduced model, and in both divisions the ordering of the three boundaries is the same. The unreduced model has the largest domain, the Craig--Bampton model comes next, and the Guyan model has the smallest. The contraction is larger in division II than in division I for both methods.

The recovery relations of Section \ref{sec:multiple-delay} \rev{help interpret the difference between the two reduced closed loops. The Guyan recovery of the physical substructure depends entirely on delayed retained coordinates. Craig--Bampton adds internal modal coordinates that do not pass through the actuators, while their reduced equations remain coupled to the delayed coordinates. Both formulations still contain two delay channels. The larger Craig--Bampton region in the reference comparison is therefore a result for the specified model--regulator pair, consistent with retaining internal modal dynamics. The recovery relation alone does not determine the size of that region.}

The stability domain of the unreduced model is itself smaller in division II than in division I. The two divisions differ in the extent of the physical substructure and in the position of the second actuator, and the combined change lowers the delay margin before any condensation is applied. Condensation adds to this loss, and the spacing between the boundaries within one panel is produced by the condensation together with the regulator supplied by each condensed model.

\FloatBarrier
\subsection{Modal redistribution ratio and applicability of the two methods}
\label{sec:energy-applicability}

The modal redistribution ratio of Eq. \eqref{eq:energy-change} is listed in Table \ref{tab:energy}. \rev{It is evaluated from the full-frame accuracy models using their first five modes. Guyan gives \(-0.0414\) in division I and \(0.8914\) in division II; the corresponding Craig--Bampton values are \(0.0039\) and \(0.0365\). Craig--Bampton is closer to zero in magnitude in both divisions, indicating a smaller change in the selected coordinate shares. The signed values do not have the same ordering in both divisions.}

\begin{table}[htbp]
	\centering
	\caption{Modal redistribution ratio of the four condensed models.}
	\label{tab:energy}
	\begin{tabular}{lcc}
		\toprule
		Division & Guyan & Craig--Bampton\\
		\midrule
		I  & \rev{-0.0414} & \rev{0.0039}\\
		II & \rev{0.8914} & \rev{0.0365}\\
		\bottomrule
	\end{tabular}
\end{table}

\rev{Figure \ref{fig:modal-shares} resolves this change into the physical-coordinate distribution and the two actuator contributions. In division II, Guyan increases the first-storey share by 113.87\% and decreases the third-storey share by 24.73\%, whose signed relative contributions sum to \(\Delta E=0.8914\). Craig--Bampton gives changes of 3.94\% and \(-0.29\%\) at the same coordinates. The figure thus locates the redistribution that the scalar index combines, while also showing why its sign cannot be read as a change in total modal energy.}

\begin{figure}[htbp]
\centering
\includegraphics[width=.98\textwidth]{fig16_modal_shares.pdf}
\caption{\rev{Low-order modal-share distribution and actuator contributions. (a,c) Shares on the fifteen physical coordinates. (b,d) Signed relative changes at the two actuated coordinates, whose sum gives the ratio in Table \ref{tab:energy}. The top and bottom rows correspond to divisions I and II, respectively.}}
\label{fig:modal-shares}
\end{figure}

\rev{The modal-share ratio is available once a reduced model has been formed and can be used alongside frequency and response errors to inspect the effect of a retained-coordinate choice. Its separate coordinate contributions are particularly useful when a horizontal coordinate is condensed. Stability is assessed from the delayed closed-loop model, including its regulator and integration algorithm; the present ratio provides no threshold for the admissible channel delays.}

The accuracy and the stability results together indicate the range of application of each method. Craig--Bampton condensation \rev{limits the first-two-mode frequency error to 0.284\% and the first-storey NRMSE to 0.138\% over the evaluated earthquake record and chirp bands}, and its stability boundary is the closer of the two to that of the unreduced model. It is therefore the method for a condensed principal coordinate and for an excitation carrying energy near or above the fundamental frequency. \rev{The full-frame Guyan models use six and five generalized coordinates, compared with nine and eight for Craig--Bampton.} In division I its second natural frequency is in error by 5.411\% and \rev{its NRMSE reaches 2.010\% in the fundamental-frequency band and 2.041\% in the highest band}, which sets the tolerance required of the application. In division II its frequency error reaches \rev{30.574\% and its NRMSE reaches 26.858\%}, more than an order of magnitude above the Craig--Bampton values, and it also gives the smallest stability domain.

The three indices do not carry equal weight in this judgement. The frequency error and the NRMSE of the Guyan model of division II are the largest values of Tables \ref{tab:modal} and \ref{tab:nrmse}, whereas the MAC of the same model remains above 0.92. A modal correlation index evaluated on the retained coordinates alone is therefore insufficient to accept a condensed model for RTHS, and it is applied together with a frequency error and a time-domain error.




\section{Conclusions} \label{sec6}

This study establishes a joint assessment of condensation accuracy and delay-dependent stability for MDOF RTHS with limited actuation channels. The common projection formulation distinguishes the approximation of structural response from the approximation of interface compatibility, and the channel-wise delay representation identifies how reconstructed responses enter the feedback loop. Guyan reconstruction depends entirely on the retained physical DOFs. Craig--Bampton reconstruction additionally includes numerically evolved fixed-interface modal DOFs, \rev{which are not directly delayed by the actuators and remain dynamically coupled to the delayed channels.} This distinction connects the reduction basis to the closed-loop response rather than treating condensation solely as a reduction of computational cost.

The numerical results for the two divisions of the three-storey frame show that Craig--Bampton condensation retains substantially greater response fidelity under the same two-channel actuation constraint. Its largest relative error among the first two natural frequencies is \rev{0.284\%}. With normalization by the full-record response range, the first-storey displacement NRMSE \rev{has a maximum of 0.138\%} for the earthquake record and all evaluated chirp bands. For the Guyan model of division II, the corresponding maxima reach \rev{30.574\% and 26.858\%}, respectively. \rev{This retained-coordinate set eliminates the middle-storey horizontal DOF. At that reconstructed coordinate, the Guyan earthquake/chirp NRMSE values are 10.798\%/11.749\%, compared with 0.031\%/0.034\% for Craig--Bampton.} The retained-DOF MAC nevertheless remains above 0.92, demonstrating that mode-shape correlation alone does not establish acceptable dynamic accuracy. Guyan condensation offers the smaller assembled model, \rev{with six/five generalized DOFs compared with nine/eight for Craig--Bampton in divisions I/II}, but the additional modal DOFs substantially improve fidelity when an important horizontal DOF is condensed.

\rev{In the reference delay-domain comparison, the Craig--Bampton model--regulator pair has a larger stable region than the Guyan pair in both divisions.} Both regions contract relative to the unreduced reference, which retains ideal, delay-free coupling at the other interface DOFs. \rev{The full-frame modal redistribution ratios are 0.0039 and 0.0365 for Craig--Bampton and \(-0.0414\) and 0.8914 for Guyan in divisions I and II, respectively. Their coordinate-wise contributions describe the change in low-order modal distribution; closed-loop poles remain the criterion for delay stability.} Thus, the smaller state dimension of Guyan condensation must be weighed against response error and admissible channel delays, whereas Craig--Bampton condensation \rev{combines higher response accuracy with the larger reference delay domains in the cases examined.} These conclusions apply to the studied linear models with independent constant channel delays, the specified regulator design and numerically represented inertia. Extension to time-varying delays or experimentally returned inertial forces requires the corresponding feedback dynamics to be included in the stability model.


\FloatBarrier
\bibliographystyle{elsarticle-num-names}
\bibliography{references}

\end{document}
